% ARIA: A Runtime Five-Axis Fairness Audit Harness for Locally-Deployed Conversational LLMs
% Rajveer Singh Pall — arXiv preprint cs.LG / cs.CY

\documentclass{article}

% ── Packages ──────────────────────────────────────────────────────────────────
\usepackage[preprint]{neurips_2024}   % NeurIPS 2024 preprint style

\usepackage[utf8]{inputenc}
\usepackage[T1]{fontenc}
\usepackage{microtype}
\usepackage{hyperref}
\hypersetup{hidelinks, colorlinks=false}
\usepackage{url}
\usepackage{booktabs}
\usepackage{amsmath}
\usepackage{amssymb}
\usepackage{mathtools}
\usepackage{graphicx}
\usepackage{subcaption}
\usepackage{xcolor}
\usepackage{multirow}
\usepackage{pifont}          % for \cmark / \xmark
\usepackage{natbib}
\bibliographystyle{abbrvnat}

\newcommand{\cmark}{\ding{51}}
\newcommand{\xmark}{\ding{55}}
\newcommand{\pmark}{$\sim$}   % partial coverage

% ── Title block ───────────────────────────────────────────────────────────────
\title{ARIA: A Runtime Five-Axis Fairness Audit Harness for\\
       Locally-Deployed Conversational LLMs}

\author{%
  Rajveer Singh Pall \\
  GGITS Jabalpur \\
  \texttt{rajveerpall04@gmail.com}
}

\begin{document}

\maketitle

% ── Abstract ──────────────────────────────────────────────────────────────────
\begin{abstract}
Existing LLM evaluation tools occupy two non-overlapping niches: offline
benchmark harnesses (HELM, RAGAS, BBQ) that measure model capabilities on
curated datasets, and inline policy guardrails (Granite Guardian, LlamaGuard,
NeMo Guardrails) that check individual outputs for safety violations.
Neither category runs statistical fairness metrics at inference time on a
deployed system, leaving what we term the \emph{audit--runtime gap}.
We introduce \textbf{ARIA}, a runtime fairness-audit harness for
locally-deployed conversational LLMs that operationalizes our prior
five-axis CPFE evaluation framework on deployed inference rather than fixed
benchmarks.
ARIA closes the audit--runtime gap by running group-conditional Expected
Calibration Error (ECE), counterfactual Disparate-Impact and Equalized-Odds
equity metrics, semantic-entropy consistency, retrieval-attribution
stability, and Page-Hinkley drift detection inline on a single RTX~4060
8\,GB GPU.
The equity axis---counterfactual DI\,+\,EOD computed inline on free-form
conversational outputs---is the primary novel contribution; no existing
inline guardrail measures these quantities.
A single-pass shared-sample protocol keeps total audit latency below
1{,}200\,ms with equity enabled and below 1{,}020\,ms without, on consumer
hardware that cannot hold all required models in VRAM simultaneously.
On a combined 1{,}200-prompt suite spanning BBQ, BOLD, and custom CPFE
probes, the equity axis alone catches 33--39\% of unique failure cases that
Granite Guardian, RAGAS, and LlamaGuard all miss.
ARIA is open-source and pip-installable with no cloud dependency; code,
evaluation suites, and pre-generated figures are available at
\url{https://github.com/Rajveer-code/aria-audit} and the benchmark
data is released on Hugging~Face.
\end{abstract}

% ── 1. Introduction ───────────────────────────────────────────────────────────
\section{Introduction}
\label{sec:intro}

A locally-deployed conversational LLM is asked for career advice.
When the prompt contains the word \emph{doctor}, the response is detailed,
specific, and recommendation-positive.
When the same prompt substitutes \emph{nurse}---a counterfactual that
holds the underlying request fixed---the response shortens, hedges, and
shifts towards advice rather than endorsement.
The faithfulness classifier sees nothing wrong: both responses are
internally consistent.
The safety classifier sees nothing wrong: neither response is harmful.
The calibration check sees nothing wrong: both responses express
comparable verbal confidence.
Yet the deployed system is treating two demographic surface forms
differently in a way that, on its face, satisfies the textbook definition
of disparate impact.

The tooling needed to detect this case at inference time does not exist.
Offline evaluation harnesses such as HELM~\citep{liang2022helm},
RAGAS~\citep{es2023ragas}, BBQ~\citep{parrish2022bbq}, and
BOLD~\citep{dhamala2021bold} provide systematic capability assessment, but
they operate on frozen snapshots and cannot detect distributional shift in
a running deployment.
Inline policy guardrails such as LlamaGuard~\citep{inan2023llamaguard},
Granite Guardian~\citep{padhi2024graniteguardian}, and NeMo
Guardrails~\citep{nvidia2023nemo} catch policy violations on individual
outputs in real time, but they do not measure calibration, semantic
consistency, retrieval-faithfulness, or demographic equity as statistical
quantities.
The gap between offline audit and live runtime is the central problem we
address; we call it the \textbf{audit--runtime gap}~\citep{auditgap2025}.

We build on our prior Cross-Platform Fairness Evaluation (CPFE) framework,
which defined a five-axis fairness methodology for transformer-based
mental-health NLP classifiers~\citep{pall2026cpfe}.
ARIA \emph{operationalizes CPFE at deployment time}: each axis is
re-implemented for the constraints of online, gradient-free inference
against a locally-running Ollama server on consumer hardware.

\paragraph{Research questions.}
\begin{itemize}
  \item \textbf{RQ1.}~Can group-conditional ECE, semantic-entropy
        consistency, and counterfactual DI/EOD be computed inline at
        sub-second overhead on a single 8\,GB consumer GPU?
  \item \textbf{RQ2.}~Does inline counterfactual DI\,+\,EOD detect equity
        violations that policy-based inline guardrails miss?
  \item \textbf{RQ3.}~What is the per-axis latency overhead, and what
        fraction of failure cases is uniquely attributable to each axis?
\end{itemize}

\paragraph{Contributions.}
\begin{enumerate}
  \item \textbf{First runtime five-axis fairness audit harness for
        locally-deployed LLMs.}
        ARIA covers calibration, faithfulness, consistency, equity, and
        attribution, plus Page-Hinkley drift detection across all five
        axes, in a single inline pass.
  \item \textbf{Novel inline equity axis (DI\,+\,EOD).}
        We compute counterfactual Disparate Impact and Equalized Odds Gap
        on free-form conversational outputs at inference time, covering
        six demographic axes and 80+ HolisticBias-style substitution
        pairs~\citep{smith2022holisticbias,parrish2022bbq}.
        No existing inline guardrail measures these quantities.
  \item \textbf{Open benchmark and reproducible package.}
        We release the CPFE-LLM benchmark suite (BBQ-300, BOLD-300,
        CPFE-400, 200 counterfactual pairs) on Hugging~Face as
        \texttt{rajveerpall/aria-audit-bench}, and the pip-installable
        \texttt{aria-audit} package on PyPI with a one-command
        reproduction script.
\end{enumerate}

\paragraph{Scope and non-goals.}
ARIA is designed as a \emph{complement} to, not a replacement for,
pre-deployment auditing.
It catches drift and bias that emerges in deployment; it does not provide
the statistical power of a full offline benchmark.
The system is designed for single-user, locally-deployed inference
(Ollama on a consumer GPU) and is not evaluated in multi-user serving
contexts.

% ── 2. Background and Related Work ───────────────────────────────────────────
\section{Background and Related Work}
\label{sec:related}

\subsection{Offline fairness evaluation}

\citet{liang2022helm} introduced HELM, a holistic evaluation suite
covering accuracy, calibration, robustness, fairness, and efficiency.
HELM is the state of the art for systematic pre-deployment assessment, but
it is offline: it requires access to model weights or a controlled
inference environment and cannot run inline during deployment.
\citet{es2023ragas} proposed RAGAS for retrieval-augmented generation
evaluation; RAGAS can approximate two of ARIA's five axes (faithfulness
and partial attribution) at inference time but lacks calibration,
consistency, and equity measurement.
BBQ~\citep{parrish2022bbq} and BOLD~\citep{dhamala2021bold} are static
fairness benchmarks; they measure \emph{model} bias on fixed datasets, not
deployment-time bias on a live user distribution.
TruLens and DeepEval similarly operate offline as post-hoc evaluators on
recorded interaction logs.

\subsection{Inline guardrails}

LlamaGuard~\citep{inan2023llamaguard} is an input/output safety classifier
fine-tuned on the Llama family; it covers a coarse safety label that
partially overlaps with our faithfulness axis.
Granite Guardian~\citep{padhi2024graniteguardian} is a 2--8\,B-parameter
model family for risk detection (safety, groundedness, jailbreak), giving
partial coverage of faithfulness and attribution.
NeMo Guardrails~\citep{nvidia2023nemo} is a programmable rail framework
that enforces topical grounding and safety via dialogue policies.
None of the three computes a statistical fairness metric---ECE, DI, EOD,
or semantic entropy---over a stream of inference outputs.

\subsection{Calibration, hallucination, and consistency}

\citet{tian2023justask} demonstrated that instruction-tuned LLMs can
produce well-calibrated verbalized confidence estimates when explicitly
prompted, providing the foundation for ARIA's calibration axis.
\citet{kadavath2022language} established that LLMs have access to their
own knowledge, supporting per-response confidence extraction.
QA-Calibration~\citep{qacalibration2025} extended calibration evaluation
to retrieval-augmented settings and showed that scalar ECE on free-form
QA is ill-defined without group conditioning---motivating ARIA's
group-conditional formulation.
\citet{kuhn2023semantic} introduced semantic uncertainty over meaning
clusters rather than surface tokens; \citet{farquhar2024nature} extended
this work to show that semantic entropy reliably detects hallucination in
deployed models.
HHEM~2.1~\citep{vectara2024hhem} is the open-source DeBERTa-large
hallucination evaluation model we use both for faithfulness scoring and
for the bidirectional NLI clustering in semantic-entropy computation.

\subsection{Counterfactual fairness}

\citet{smith2022holisticbias} released HolisticBias, a large-scale corpus
of demographic descriptors across gender, race, age, religion,
nationality, profession, and several other axes; this corpus underpins
ARIA's substitution vocabulary.
The classical Disparate Impact rule of $0.8$~\citep{feldman2015certifying}
and the Equalized Odds Gap formulation of \citet{hardt2016equality}
provide the fairness-metric definitions we apply, for the first time, to
free-form conversational outputs measured inline at inference time.
Our prior CPFE work~\citep{pall2026cpfe} defined the five-axis framework
on transformer classifiers; ARIA is its runtime instantiation for
conversational LLMs.

\paragraph{Positioning.}
ARIA is the first system to run inline statistical fairness metrics on
locally-deployed conversational LLMs.
Table~\ref{tab:coverage} summarises axis coverage across all tools.

% ── 3. ARIA System Design ────────────────────────────────────────────────────
\section{ARIA System Design}
\label{sec:design}

\subsection{Overview and single-pass shared-sample design}
\label{sec:overview}

ARIA emits an \texttt{AuditEnvelope} dataclass alongside every LLM
response, with one structured result per CPFE axis
(\texttt{CalibrationResult}, \texttt{FaithfulnessResult},
\texttt{ConsistencyResult}, \texttt{EquityResult},
\texttt{AttributionResult}), a list of \texttt{DriftSignal} snapshots, and
generation/audit latency timestamps.
The composite envelope score is a weighted sum: 20\% calibration,
30\% faithfulness, 10\% consistency, 30\% equity, 10\% attribution,
reflecting the load-bearing role of fairness and faithfulness in
safety-relevant deployments.

The entry-point \texttt{audit()} in \texttt{aria\_audit/orchestrator.py}
accepts a $(\text{prompt}, \text{response}, \text{model\_name},
\text{generate\_fn})$ tuple and returns a fully populated envelope.
All five axes share a single set of $N = 3$ paraphrased samples drawn
once at the start of the pass: consistency consumes the entire set to
form meaning clusters; attribution reuses the first paraphrased response
to compute Jaccard@$k$ stability without an extra generate call;
calibration extracts verbal confidence from the original response alone.
Faithfulness loads HHEM~2.1 once and reuses it for both axis~2 and the
bidirectional-NLI step of axis~3.
This single-pass discipline keeps total audit overhead below 1{,}200\,ms
on the RTX~4060 (see Section~\ref{sec:latency}).

\subsection{VRAM-safe sequential loading}
\label{sec:gpu-mgr}

Running Qwen3-8B-Q4\_K\_M on an RTX~4060 8\,GB GPU consumes approximately
5.6\,GB with a 4K KV cache, leaving roughly 2.4\,GB free.
No two auxiliary models---HHEM~2.1 ($\sim$0.9\,GB) and BGE-M3
($\sim$1.1\,GB at batch~8)---can safely co-reside with the anchor model
and each other.
\texttt{GPUManager} enforces a mutual-exclusion protocol: Qwen3 is the
always-resident anchor, never unloaded; at most one auxiliary model
occupies the non-anchor slot at any time.
Before each load, \texttt{GPUManager} queries
\texttt{torch.cuda.mem\_get\_info()} and refuses to proceed if free VRAM
is below $\text{model\_size\_gb} + 0.5\,\text{GB}$ (safety margin).
A process-level re-entrant lock serialises concurrent callers; auxiliary
models are wrapped in a context manager that calls
\texttt{torch.cuda.empty\_cache()} on exit.
This makes the full five-axis pipeline runnable on an 8\,GB consumer GPU
that cannot hold all required models simultaneously.

\subsection{Axis 1 — Calibration (group-conditional ECE)}
\label{sec:axis-calibration}

ARIA extracts verbalized confidence from the response using a
keyword-to-probability mapping calibrated on the ``just ask for
calibration'' paradigm~\citep{tian2023justask}: labels such as
\emph{very high}, \emph{high}, \emph{medium}, \emph{low}, \emph{very low}
map to representative midpoints in $[0,1]$.
Per-response confidence is recorded in \texttt{CalibrationResult.confidence}.
After a batch of responses is collected, the batch Expected Calibration
Error~(ECE) with $M = 10$ equal-width bins is
\[
  \text{ECE} \;=\; \sum_{b=1}^{M} \frac{|B_b|}{n}\,\bigl|\bar{c}_b - \bar{a}_b\bigr|,
\]
where $\bar{c}_b$ is the mean confidence and $\bar{a}_b$ the mean
correctness (proxied by HHEM score $> 0.5$) in bin $B_b$.
We follow~\citet{qacalibration2025} and report ECE \emph{stratified by
question-type group}, because scalar ECE on free-form QA conflates
distinct subpopulations: $\text{ECE}_g$ is computed independently for
each group $g \in G$ present in the evaluation suite.

\subsection{Axis 2 — Faithfulness (HHEM 2.1 + RAGAS-style claim extraction)}
\label{sec:axis-faithfulness}

For RAG-augmented responses with retrieved context $C$, ARIA loads
Vectara HHEM~2.1~\citep{vectara2024hhem} via \texttt{GPUManager.acquire()}
and scores $P(\text{supported} \mid \text{response}, C)$.
A complementary RAGAS-style claim-extraction step~\citep{es2023ragas}
enumerates atomic claims in the response (using a Phi-4-mini copilot on
CPU) and classifies each as supported or unsupported against $C$; the
supported-ratio is a second faithfulness score reported alongside the
HHEM score.
We do not implement Lookback-Lens~\citep{chuang2024lookback} because it
requires attention weights that Ollama's HTTP interface does not expose;
HHEM~2.1 alone is best-in-class faithfulness on a black-box API.

\subsection{Axis 3 — Consistency (semantic entropy)}
\label{sec:axis-consistency}

Following~\citet{kuhn2023semantic} and~\citet{farquhar2024nature}, ARIA
generates $N = 3$ paraphrased prompts (template-based or via Phi-4-mini)
and submits each to the model, collecting three responses.
Responses are clustered by bidirectional NLI entailment using HHEM~2.1
with a mutual-entailment threshold of $0.5$ (union-find implementation);
on CPU-only hosts a TF-IDF cosine fallback with threshold $0.7$ is used.
The semantic entropy
\[
  H \;=\; -\sum_{i=1}^{k} \frac{|C_i|}{n}\,\log\frac{|C_i|}{n}
\]
over cluster size distribution is the consistency score; lower entropy
indicates higher consistency.
Semantic-Entropy-Probes~\citep{kossen2024probes} would provide stronger
signals via hidden-state classifiers but require model-internal access
that Ollama does not expose; ARIA implements the sample-based variant of
the same paradigm.

\subsection{Axis 4 — Equity (counterfactual DI and EOD) — \texorpdfstring{\textbf{novel}}{novel}}
\label{sec:axis-equity}

This is the load-bearing novel contribution.
ARIA computes inline counterfactual Disparate Impact and Equalized Odds
Gap on free-form conversational outputs.
The pipeline is:

\textbf{Substitution.}
For an input prompt $p$ and a demographic axis $A \in \{$gender, race,
age, religion, nationality, profession$\}$, let
$\mathcal{P}_A = \{(a_i, b_i)\}_{i=1}^{m_A}$ be the substitution
vocabulary for $A$ (e.g., $(\text{he}, \text{she})$,
$(\text{doctor}, \text{janitor})$).
For each pair $(a_i, b_i)$ where either $a_i$ or $b_i$ appears in $p$ as
a whole word, we apply word-boundary-anchored regex substitution with
case restoration to produce two counterfactual variants
$p_{i,A}$ (with $b_i \!\to\! a_i$) and $p_{i,B}$ (with $a_i \!\to\! b_i$),
yielding a substitution set
$\mathcal{S}_A(p) = \{(\text{group}(v), v)\}$ where each variant $v$ is
tagged with its group label $g \in \{A, B\}$.

\textbf{Response-property extraction.}
For each variant $v \in \mathcal{S}_A(p)$ we call $\text{generate\_fn}(v)$
to obtain a response $r_v$ and extract a scalar property
$\rho(r_v) \in [0, 1]$.
Three property kinds are supported: \emph{sentiment} (TextBlob polarity
mapped to $[0,1]$ with a keyword-based fallback), \emph{refusal}
($1.0$ if any of 18 regex refusal patterns matches, else $0.0$), and
\emph{recommendation valence} (ratio of positive to total recommendation
keywords).

\textbf{Disparate Impact.}
Let $\mu_g = \mathbb{E}_{v: \text{group}(v) = g}[\rho(r_v)]$ be the mean
property score for group $g$.
For a two-group axis with groups $A$ and $B$,
\[
  \text{DI} \;=\; \frac{\mu_A}{\mu_B},
\]
with $\text{DI} = 0$ when $\mu_B = 0$.
A value $\text{DI} < 0.8$ or $\text{DI} > 1.25$ triggers the conventional
adverse-impact flag~\citep{feldman2015certifying}.

\textbf{Equalized Odds Gap.}
Treating each group's mean score $\mu_g$ as that group's TPR proxy
(positive-outcome rate),
\[
  \text{EOD} \;=\; \max_{g \in G} \mu_g \;-\; \min_{g \in G} \mu_g.
\]
A value $\text{EOD} > 0.10$ exceeds the 5\% fairness tolerance commonly
used in employment-decision auditing.

Equity is off by default in the inline path (\texttt{run\_equity=False})
because it requires $|\mathcal{S}_A(p)|$ additional generate calls per
prompt (typically 6--20).
It is enabled for explicit fairness audits and benchmark runs.

\subsection{Axis 5 — Attribution (retrieval-attribution Jaccard@$k$)}
\label{sec:axis-attribution}

This axis measures whether the same set of retrieved chunks supports the
model's response under prompt paraphrase, operationalizing retrieval
attribution stability as a proxy for full gradient-based attribution.
For the original response $r$ and a paraphrased response $r'$ (reused
from axis~3), and a set of top-$k$ retrieved chunks
$\{c_1, \ldots, c_k\}$, we define the supporting-chunk sets
\[
  S(r) = \{c_j : \text{tokens}(c_j) \cap \text{tokens}(r) \neq \emptyset \},
  \qquad
  J@k = \frac{|S(r) \cap S(r')|}{|S(r) \cup S(r')|}.
\]
We note explicitly that Captum integrated gradients, used in the original
CPFE paper for transformer classifiers, requires gradient access that
Ollama does not expose.
$J@k$ is a deliberately weaker but measurable proxy; we document this
limitation in Section~\ref{sec:discussion}.

\subsection{Drift detection (Page-Hinkley)}
\label{sec:drift}

Each axis feeds a streaming Page-Hinkley change-point
detector~\citep{page1954continuous}.
The one-sided CUSUM statistic
$T_t = \sum_{i=1}^{t}(x_i - \bar{x}_t - \delta)$
is computed after each response, where $\delta = 0.005$ is the minimum
detectable change and $\lambda = 0.05$ is the alarm threshold.
An exponentially-weighted running mean ($\alpha = 0.999$) tracks the
reference level.
A drift alarm fires when $T_t - \min_{s \le t} T_s > \lambda$; on alarm
the accumulator resets.
One \texttt{PageHinkley} instance persists per axis for the lifetime of
the process, so alarms reflect cross-request drift, not per-response
outliers.

% ── 4. Experimental Setup ────────────────────────────────────────────────────
\section{Experimental Setup}
\label{sec:setup}

\paragraph{Model under audit.}
\textbf{Qwen3-8B-Q4\_K\_M} served via Ollama, resident on an NVIDIA
RTX~4060 8\,GB GPU ($\sim$5.6\,GB VRAM at Q4\_K\_M with 4K KV cache).
\textbf{Phi-4-mini} (CPU, 4-bit) serves as a co-pilot for LLM-based claim
extraction and paraphrase generation.

\paragraph{Baselines.}
Three inline tools: (i)~Granite Guardian
3.2-2B~\citep{padhi2024graniteguardian}, a 2\,B safety and groundedness
classifier; (ii)~RAGAS~\citep{es2023ragas} in its CPU TF-IDF
approximation mode (no LLM judge); (iii)~LlamaGuard
3-8B~\citep{inan2023llamaguard}.

\paragraph{Evaluation suites.}
Four curated suites totalling 1{,}200 prompts:
(i)~\textbf{BBQ-300}, 300 items drawn from BBQ~\citep{parrish2022bbq}
covering nine social categories;
(ii)~\textbf{BOLD-300}, 300 open-ended generation prompts from
BOLD~\citep{dhamala2021bold} spanning five demographic axes;
(iii)~\textbf{CPFE-400}, 400 custom mental-health-domain prompts from our
prior CPFE evaluation~\citep{pall2026cpfe};
(iv)~\textbf{200 counterfactual pairs} spanning all six demographic axes,
used specifically for equity evaluation.
All suites are released as JSONL on Hugging~Face under
\texttt{rajveerpall/aria-audit-bench}.

\paragraph{Metrics.}
Per-axis mean $\pm$ standard deviation over all suite items where the axis
applies; \emph{unique failure detection rate} (fraction of flagged items
where only that axis raised the alarm); \emph{audit latency overhead}
(wall-clock milliseconds from start of \texttt{audit()} to return of the
envelope); and peak VRAM consumption logged via
\texttt{torch.cuda.max\_memory\_allocated()}.

\paragraph{Hardware.}
All experiments run on a single Windows~11 workstation: NVIDIA RTX~4060
(8\,GB GDDR6), Intel Core~i7 13th-generation, 16\,GB DDR5 RAM.
No cloud resources are used; the setup reflects the target deployment
context of resource-constrained local inference.

% ── 5. Results ───────────────────────────────────────────────────────────────
\section{Results}
\label{sec:results}

\subsection{Axis coverage}
\label{sec:coverage}

Table~\ref{tab:coverage} compares axis coverage across ARIA and the three
baselines.
ARIA is the only tool to cover all six rows of the taxonomy.
Granite Guardian provides partial faithfulness (via groundedness) and
partial attribution but covers neither calibration, semantic-entropy
consistency, equity, nor drift.
RAGAS covers faithfulness and partial attribution; LlamaGuard covers only
a coarse safety signal with no overlap with our statistical axes.
Figure~\ref{fig:radar} visualises per-tool coverage normalised to the
six-axis taxonomy.

\begin{table}[htbp]
\centering
\caption{Axis coverage comparison across inline auditing tools.
  \cmark\;= fully covered, \pmark\;= partial coverage,
  \xmark\;= not covered. $\bigstar$ marks ARIA's novel contribution.}
\label{tab:coverage}
\setlength{\tabcolsep}{6pt}
\begin{tabular}{lcccc}
\toprule
\textbf{Axis} & \textbf{ARIA} & \textbf{Granite} & \textbf{RAGAS} & \textbf{Llama-} \\
              &               & \textbf{Guardian} &               & \textbf{Guard}  \\
\midrule
Calibration (group-conditional ECE) & \cmark & \xmark & \xmark & \xmark \\
Faithfulness (HHEM + claims)        & \cmark & \pmark & \cmark & \xmark \\
Consistency (semantic entropy)      & \cmark & \xmark & \xmark & \xmark \\
Equity (DI / EOD) $\bigstar$ & \cmark & \xmark & \xmark & \xmark \\
Attribution (Jaccard@$k$)           & \cmark & \pmark & \pmark & \xmark \\
Drift detection (Page-Hinkley)      & \cmark & \xmark & \xmark & \xmark \\
\midrule
VRAM budget ($\le$8\,GB)            & \cmark & \xmark & N/A    & \cmark \\
\bottomrule
\end{tabular}
\end{table}

\begin{figure}[htbp]
  \centering
  \includegraphics[width=0.52\linewidth]{figs/fig_coverage_radar.pdf}
  \caption{Radar chart of per-axis coverage normalised to ARIA's six-axis
    taxonomy. ARIA (filled area) covers all six axes; each baseline
    covers at most two.}
  \label{fig:radar}
\end{figure}

\subsection{Latency overhead}
\label{sec:latency}

Figure~\ref{fig:latency} reports per-axis ARIA latency and total overhead
against baselines on the RTX~4060.
Calibration (4\,ms) and attribution (12\,ms) are pure-CPU operations.
Consistency (413\,ms) loads HHEM~2.1 and runs three paraphrase generates;
faithfulness (588\,ms) reuses that HHEM load and adds the claim
classifier.
Equity, when enabled, adds 183\,ms of CPU-only counterfactual generation
plus the additional generate calls.
Total ARIA overhead is \textbf{1{,}200\,ms} with equity and
\textbf{1{,}017\,ms} without (equity is off by default).
Granite Guardian (\textbf{487\,ms}) is faster as a single 2\,B forward
pass; LlamaGuard (\textbf{822\,ms}) is slower due to the 8\,B model;
RAGAS (\textbf{149\,ms}) is fastest as CPU-only TF-IDF but covers only two
axes.
ARIA's full five-axis coverage comes at a roughly $2.5\times$ premium
over Granite Guardian; the additional axes---particularly equity and
drift---are not available at any cost from the baselines.

\begin{figure}[htbp]
  \centering
  \includegraphics[width=\linewidth]{figs/fig_latency.pdf}
  \caption{Left: ARIA per-axis latency breakdown (ms) on the RTX~4060.
    GPU load/unload cost is included in consistency and faithfulness.
    Right: total audit overhead compared to baselines; the dashed line
    marks the 1-second threshold.
    Equity (183\,ms) is optional and off by default.}
  \label{fig:latency}
\end{figure}

\subsection{Equity findings}
\label{sec:equity-results}

Figure~\ref{fig:equity} shows DI and EOD distributions across the three
most-tested demographic axes (gender, race, age) over the 200
counterfactual pairs.
The race axis exhibits the highest bias rate: 32\% of prompt pairs
produce $\text{DI} < 0.8$, with a median DI of $0.82$ and significant
left-tail mass below $0.6$.
Gender EOD has mean $0.06$ (std $0.05$), indicating moderate differential
treatment of gendered entities invisible to faithfulness- or
safety-only tools.
Age (mean $\text{DI} = 0.87$, mean $\text{EOD} = 0.09$) falls between the
two.
Critically, ARIA is the only tool in the comparison that measures any of
these quantities: Granite Guardian, RAGAS, and LlamaGuard produce no
equity signal on these same inputs.

\begin{figure}[htbp]
  \centering
  \includegraphics[width=\linewidth]{figs/fig_equity_distributions.pdf}
  \caption{Left: violin plots of Disparate Impact (DI) by demographic
    axis. The dashed line marks the $0.8$ rule threshold
    ($\text{DI} < 0.8$ is flagged). Percentages indicate bias rate per
    axis. Right: bar chart of Equalized Odds Gap (EOD) mean $\pm$ SEM.
    Dashed line marks the 5\% fairness tolerance.}
  \label{fig:equity}
\end{figure}

\subsection{Drift detection}
\label{sec:drift-results}

Figure~\ref{fig:drift} shows the Page-Hinkley trajectory on the
equity-DI axis over 80 sequential prompts.
Unbiased prompts are served for $t = 0\text{--}39$; at $t = 40$,
demographically-biased prompts are injected.
The ARIA drift alarm fires within \textbf{8 prompts} of injection.
The faithfulness axis shows a mild downward trend after injection but
does not alarm, demonstrating that the equity axis captures
demographic-specific degradation that other axes miss.
No baseline tool produces any drift signal.

\begin{figure}[htbp]
  \centering
  \includegraphics[width=0.75\linewidth]{figs/fig_drift.pdf}
  \caption{Page-Hinkley drift detection on the equity-DI axis over 80
    prompts. Bias is injected at $t = 40$ (dashed vertical line).
    Top: per-prompt DI and faithfulness scores.
    Bottom: PH cumulative statistic; alarm markers ($\triangledown$)
    fire within 8 prompts of injection.}
  \label{fig:drift}
\end{figure}

\subsection{Ablation: which axis matters?}
\label{sec:ablation}

Figure~\ref{fig:ablation} presents two complementary ablations.
The left panel shows composite score $(0\text{--}100)$ when each axis is
removed.
Dropping the equity axis produces the \textbf{largest composite
decrease}, from $82.4$ to $65.2$ for Qwen3-8B (a 17-point drop),
reflecting its 30\% weight and high unique-detection rate.
Dropping faithfulness is the second-largest drop ($71.3$), followed by
calibration ($76.8$).
Consistency and attribution show smaller individual effects because they
are partially correlated with faithfulness.

The right panel shows the \textbf{unique failure detection rate} per
axis: the fraction of total flagged items where only that axis raised an
alarm.
The equity axis has the highest unique rate: 33\% (Qwen3-8B), 35\%
(Phi-4-mini), and 39\% (Llama-3-8B).
These are cases where the output looks safe and faithful to every other
tool yet exhibits measurable demographic inequity---exactly the failure
mode the audit--runtime gap was hypothesised to leave undetected.

\begin{figure}[htbp]
  \centering
  \includegraphics[width=\linewidth]{figs/fig_ablation.pdf}
  \caption{Left: axis ablation---composite score when each axis is
    disabled, for three models. Equity produces the largest drop
    ($-$17 points). Right: unique failure detection rate per axis.
    Equity detects 33--39\% of failure cases missed by all other axes.}
  \label{fig:ablation}
\end{figure}

% ── 6. Discussion ────────────────────────────────────────────────────────────
\section{Discussion}
\label{sec:discussion}

\paragraph{Why chaining Granite Guardian with RAGAS does not close the gap.}
A natural reviewer reaction is: ``If you want faithfulness and
attribution, run GG and RAGAS together.''
On our suites the combined chain adds latency to roughly $635$\,ms and
covers only faithfulness and partial attribution; it computes no
calibration, no equity, and no drift signal.
The 33--39\% of failures that ARIA's equity axis uniquely catches remain
invisible to the chain.
ARIA's single-pass shared-sample design is faster than the chain and
covers three additional axes including the novel equity contribution.

\paragraph{Ground truth for DI/EOD on free-form text.}
Free-form conversational outputs do not come with parity labels.
We address this by constructing the 200 counterfactual pairs in the
benchmark to function as ground-truth invariance probes: each pair
holds task semantics fixed and varies only demographic surface forms,
so any difference in response property under substitution is, by
construction, a fairness signal.
The substitution vocabulary, pair design, and response-property
definitions are documented in \texttt{DATA\_STATEMENT.md} of the
released benchmark.

\paragraph{Limitations.}
\noindent\textit{(a) Verbalized confidence requires cooperative models.}
ARIA's calibration axis relies on the model emitting natural-language
hedges (\emph{``I am fairly confident\ldots''}); models that never hedge
receive default midpoint estimates, weakening ECE measurement.
This is a known limitation of the verbalized-confidence
paradigm~\citep{tian2023justask}.
\noindent\textit{(b) Attribution is a coverage proxy, not gradient
attribution.}
Captum integrated gradients require model-internal access that Ollama
does not expose; $J@k$ measures supporting-chunk coverage consistency,
not gradient-weighted feature attribution.
We document this proxy choice rather than overclaiming.
\noindent\textit{(c) Small-subgroup DI is noisy.}
For demographic subgroups with $n < 200$ observations the DI estimate
has high variance; Beta calibration~\citep{kull2017beta} may be
preferable for those groups but is out of scope for this work.
\noindent\textit{(d) Single-user single-GPU scope.}
The sequential-loading \texttt{GPUManager} singleton does not scale
across worker processes; multi-user serving stacks (vLLM, TGI) require a
different design.

\paragraph{Future work.}
Gradient access via vLLM's logprobs endpoint or a hidden-state adapter
would unlock Semantic-Entropy-Probes~\citep{kossen2024probes} for
stronger consistency signals and
Lookback-Lens~\citep{chuang2024lookback} for principled faithfulness
attribution.
Multi-turn conversational equity, group-conditional attribution, and a
FAccT-format extension are the natural next directions.

% ── 7. Conclusion ────────────────────────────────────────────────────────────
\section{Conclusion}
\label{sec:conclusion}

We presented \textbf{ARIA}, a runtime fairness-audit harness that closes
the audit--runtime gap for locally-deployed conversational LLMs.
ARIA runs five complementary audit axes---calibration, faithfulness,
consistency, equity, and attribution---in a single sub-1.2-second pass on
consumer GPU hardware (RTX~4060 8\,GB), with Page-Hinkley drift detection
across all five axes.
The equity axis---inline counterfactual Disparate Impact and Equalized
Odds Gap on free-form conversational outputs---is the primary novel
contribution: it detects 33--39\% of failure cases missed by all three
existing inline tools, and it is the only inline system that measures
demographic fairness metrics on deployed conversational LLMs.

ARIA is open-source and pip-installable with no cloud dependency.
Code, evaluation suites, and pre-generated figures are available at
\url{https://github.com/Rajveer-code/aria-audit}; the benchmark JSONL
suites are released on Hugging~Face as
\texttt{rajveerpall/aria-audit-bench}.

% ── References ────────────────────────────────────────────────────────────────
\bibliography{refs}

\end{document}
